Coordinated rhythmic limb movement in mammals is controlled primarily by
spinal cord circuitry~\cite{grahamBrown1911,kiehn2016}. Locomotion is one of
the most fundamental motor functions of the nervous system: beyond mobility
itself, walking underpins independence, social participation, and quality of
life. The loss of locomotor function after spinal cord injury (SCI) is
consequently one of the most disabling neurological impairments, affecting
millions of people worldwide. Substantial advances in rehabilitation and
neuromodulation have improved outcomes for some patients, but restoring
coordinated locomotion after SCI remains a major clinical challenge, in part
because how residual spinal circuits reorganize and adapt to injury, sensory
experience, and rehabilitation is still incompletely understood.

For decades, experimental and computational work has shown that the spinal
cord contains intrinsic neuronal networks capable of generating rhythmic
motor patterns without continuous supraspinal control. These spinal central
pattern generators (sCPGs) produce the alternating flexor--extensor bursting
that characterizes locomotion, a pattern preserved across
species~\cite{rybak2006,kiehn2016}. Locomotor-like activity persists even
after extensive disruption of descending pathways, underscoring the autonomy
of spinal circuits and their potential contribution to recovery after
injury. Locomotion in intact animals, however, is not generated by an
isolated spinal oscillator: it emerges from continuous interaction among
descending commands, spinal interneuronal circuits, peripheral sensory
feedback, and the biomechanics of the musculoskeletal system.

Among these interacting mechanisms, sensory feedback plays a particularly
important role. Group Ia and Ib afferents from muscle spindles and Golgi
tendon organs continuously inform the spinal cord about limb position,
muscle length, and loading, while cutaneous afferents signal contact with
the ground. These inputs do not merely modulate locomotion; they actively
shape the timing, stability, and adaptability of the locomotor rhythm.
Seminal work by Pearson and Rossignol showed that proprioceptive feedback
contributes to phase transitions, load compensation, and inter-joint
coordination during walking, establishing locomotion as a closed-loop
sensorimotor process rather than a purely centrally generated
rhythm~\cite{pearson1995,rossignol2006}. After SCI, sensory feedback takes on
even greater importance: it is a route through which residual spinal
circuits can still be recruited and reorganized once descending drive is
reduced.

The capacity of spinal networks to adapt to sensory experience ultimately
rests on synaptic plasticity. Activity-dependent plasticity lets the nervous
system modify synaptic strength in response to patterns of neuronal
activity, allowing motor circuits to reorganize after injury and in response
to rehabilitation. Rehabilitation strategies such as body-weight-support
treadmill training (BWSTT), robotic-assisted gait therapy, and epidural
spinal cord stimulation (ESCS) are thought to derive much of their
effectiveness from engaging and guiding plastic change within spinal
sensorimotor networks, rather than from regeneration of damaged descending
pathways alone~\cite{courtine2009,edgerton2008}. The specific mechanisms by
which sensory signals drive this reorganization, however, remain
incompletely understood.

Spike-timing-dependent plasticity (STDP) is a biologically plausible
candidate mechanism. STDP modifies synaptic efficacy according to the
relative timing of pre- and postsynaptic activity, letting neural circuits
self-organize on the basis of experience~\cite{biPoo1998,morrison2007}. It
has been studied extensively in cortical and hippocampal systems, but its
role within spinal locomotor networks remains largely unexplored. In
principle, STDP offers an attractive account of how locomotor circuits could
autonomously tune sensorimotor connectivity to improve stability,
efficiency, and coordination during a repetitive movement like walking;
whether such self-organizing principles actually operate in spinal locomotor
circuits bears on both basic neurobiology and rehabilitation neuroscience.

Computational modeling makes it possible to manipulate individual
mechanisms in ways that are not feasible in vivo. A lineage of increasingly
detailed spinal sCPG models --- Rybak, Shevtsova, Danner, Zhang and
colleagues --- has progressively refined the architecture of the mammalian
locomotor circuit, using Hodgkin--Huxley dynamics with persistent-sodium-driven
intrinsic bursting and genetically identified interneuron populations (V0V,
V0D, V1, V2a, V3, Shox2) to reproduce left--right coordination, interlimb
gaits, and speed-dependent gait
expression~\cite{rybak2006,shevtsova2015,danner2017,zhang2022}. These models,
however, rely on fixed, hand-curated synaptic weights, and their most recent
iteration explicitly names its own simplifications: no motoneurons, no
biomechanics or sensory feedback from the limbs, and no Ia/Ib reflex
circuitry~\cite{zhang2022}. Existing sCPG models, in short, cannot address
how plasticity and sensory feedback interact to shape locomotor behavior
over time, because neither is present in the model.

The present study addresses that gap. We extend this computational sCPG
lineage with four mechanisms: (1) STDP-driven self-organization of the
descending (BS$\to$RG) and cutaneous (CUT$\to$RG) synapses, so afferent
gains are discovered by the network rather than hand-tuned; (2) closed-loop,
force- and length-driven Ia feedback projecting back into the
reciprocal-inhibition interneurons (InE/InF) that gate the rhythm generators
(RG-E/RG-F); (3) explicit motoneuron pools (M-E/M-F) and Hill-style muscle
proxies that turn rhythm-generator output into graded force and length; and
(4) an externally paced, sequential Ia activation reproducing the
heel$\to$mid$\to$toe pressure transfer of stance. Rather than prescribing
the strength of the key sensorimotor pathways, we let them self-organize
through ongoing activity, and we ask how the resulting locomotor pattern
degrades as sensory feedback is withdrawn while a reduced, tonic
descending drive is preserved throughout --- the model's closest analogue
to motor-incomplete SCI, and the central manipulation in locomotor
rehabilitation for that population.

Beyond its computational contribution, this model has translational
motivation. People with SCI often recover a degree of stepping when sensory
input is augmented, through BWSTT or ESCS, but the neural mechanisms behind
that improvement are still debated. By combining sensorimotor feedback,
spinal plasticity, and locomotor output in one closed-loop computational
framework, the present model offers a mechanistic platform for asking how
rehabilitation paradigms might exploit residual spinal plasticity to restore
movement, a step toward the kind of patient-specific digital model that
could eventually help optimize rehabilitation after neurological injury.

Section~\ref{sec:acronyms} lists the acronyms used throughout the paper.
Section~\ref{sec:methods} describes the model architecture, the Izhikevich
and Hodgkin--Huxley neuron implementations, the STDP rule, the sensory
pathways, and the paced-gait protocol. Section~\ref{sec:results} reports how
the plastic synapses converge to a shared set-point independent of
initialization (\autoref{sec:results:weights}), how that convergence tracks
rehabilitation-relevant force output (\autoref{sec:results:force}), how
network activity compares across the $5\times5$ grid of locomotion modes and
STDP rates (\autoref{sec:results:comparison}), and how a constant cutaneous
drive rescues stepping under simulated unloading, an epidural-stimulation
analogue (\autoref{sec:results:epidural}). Section~\ref{sec:discussion}
relates these findings to the Rybak/Shevtsova/Danner/Zhang lineage and to the
experimental SCI-rehabilitation literature, and discusses limitations and
future extensions.
